\chapter{Methodology}
\section{Overview and Model Architecture}

The model is a geospatial techno-economic assessment framework that calculates the delivered cost of green hydrogen from approximately 6{,}000 global production locations to five European import gateways. All costs are expressed in EUR/kg H\textsubscript{2}. The target year for all results is 2040.

The framework is structured as a sequential pipeline:
\begin{enumerate}
    \item Spatial data loading and cost of capital assignment
    \item Generation cost computation
    \item Transport cost computation and carrier selection
    \item Country capacity allocation
    \item Corridor aggregation and supply curve construction
    \item Monte Carlo uncertainty and portfolio analysis
    \item Multi-objective dispatch and scenario matrix
    \item Repeat monte carlo analysis
\end{enumerate}

The spatial grid, routing architecture, and core LCOH formulation build on the framework established by \cite{Collis2022}. This includes the generation cost and transport cost calculations. The present model extends that foundation in six respects: (1) country and technology specific WACC; (2) multi-objective dispatch; (3) the European Hydrogen Backbone (EHB) corridor framework; (4) IEA projects database for capacity allocation; (5) three demand scenario structures; and (6) portfolio optimisation tool.

\section{Spatial Grid and Input Data}

\subsection{Onshore Grid}

The primary spatial input contains 5{,}970 onshore grid points at approximately 1-degree resolution across all continents. Each point contains:

\begin{itemize}
    \item Longitude, Latitude
    \item Solar Energy Potential [kWh/kWp/yr] --- location-specific photovoltaic yield from \cite{GSA}
    \item Wind Power Density [W/m\textsuperscript{2}] --- annual mean wind power density at hub height from \cite{GWA}
    \item Country, ISO\_A3, Continent, Region
    \item Nearest port assignment and overland road distances (pre-computed via the OSRM Open Source Routing Machine API, see \S\ref{sec:transport})
\end{itemize}

\subsection{Technology Cost Trajectories}

Annual learning curves for renewable and electrolyser capital costs are derived from the Danish Energy Agency, and the IEA Global Hydrogen Review 2025, IRENA, and DNV Energy Transition Outlook \citep{IEA2026, DNV2025, DEA2026, IRENA2022}. All cost curves are parameterised as years elapsed since 2026, capped at 24 (corresponding to 2050). The full cost curve equations are set out in \S\ref{sec:gen-cost}.

\subsection{IEA Hydrogen Projects Database}

The IEA Hydrogen Production Projects Database contains 2{,}662 green hydrogen projects \citep{IEA2026}. Each record includes: project name, country, technology, status, capacity [kt H\textsubscript{2}/yr], expected online date, and end use. This database is used to derive near term deployment anchors for 2030 in the capacity allocation model (\S\ref{sec:capacity}). Project records are filtered by status and weighted as follows:

\begin{table}[htbp]
    \centering
    \caption{IEA project database status weighting.}
    \label{tab:iea-weights}
    \begin{tabular}{lc}
        \toprule
        Project Status & Weight \\
        \midrule
        Operational & 100\% \\
        FID / FEED  & 100\% \\
        Feasibility study & 30\% \\
        Concept & 0\% \\
        \bottomrule
    \end{tabular}
\end{table}

\subsection{Cost of Capital}
\label{subsec:wacc}

The model applies country and technology specific weighted average cost of capital (WACC) values, rather than a single global discount rate that is used in \cite{Collis2022}. Two separate WACC series are used per country:

\begin{itemize}
    \item \textbf{Renewable electricity WACC}: derived from the Damodaran country risk methodology applied to onshore wind and solar projects, covering 177 countries. This reflects sovereign risk, currency risk, and the maturity of the local renewable energy market.
    \item \textbf{Electrolyser WACC}: derived from the same Damodaran framework applied to electrolysis projects.
\end{itemize}

Both series are sourced from \cite{Hatton2025}. A regional WACC look-up table is maintained as a fallback for the approximately 17 countries not covered by the Hatton datasets. Where country-level WACC data are available, they take precedence.

These WACC values are applied via the Capital Recovery Factor (CRF) in all CAPEX annualisation calculations.

\subsection{Socio-Environmental Auxiliary Data}

Two auxiliary datasets are used in the multi-objective portfolio analyses (\S\ref{sec:mc-multi} and \S\ref{sec:strategic-dispatch}):

\begin{itemize}
    \item \textbf{Governance / Security:} Transparency International Corruption Perceptions Index (CPI), 2025 edition. Scores are on a 0--100 scale (0 = highly corrupt, 100 = very clean). Used as a country-level governance proxy: a country with a high CPI score is preferred as a lower-risk supply partner \citep{CPI2025}.
    \item \textbf{Water availability:} WRI Aqueduct 4.0 Baseline Water Stress (BWS) index. BWS scores are on a 0--5 scale (0 = low stress, 5 = extremely high stress). Green hydrogen production via electrolysis requires water input; high water stress in a source country represents a supply chain risk and is penalised in the portfolio objective function \citep{wri2023}. 
\end{itemize}

\section{Generation Cost Model}
\label{sec:gen-cost}

\subsection{Overview}

The generation cost model from \cite{Collis2022} computes the levelised cost of hydrogen at the production site [EUR/kg H\textsubscript{2}] for every grid point simultaneously. It processes all locations together in one batch operation, which is faster and more efficient. The total generation cost is:

\begin{equation}
    \text{LCOH}_{\text{gen}} = \frac{\text{Annualised cost of electricity supply} + \text{Annualised cost of electrolyser system}}{\text{Annual H}_2\text{ production}}
\end{equation}

Because costs are evaluated on a per-kg basis, the absolute production volume cancels out and does not need to be specified and only specific energy and power requirements enter the calculation.

\subsection{Technology Cost Curves}

All costs are computed as a function of $\text{year\_diff} = \max(0, \min(24, \text{year} - 2026))$.

\paragraph{Solar PV CAPEX:}
\begin{align}
    \text{capex}_{\text{solar}} &= 700 \times 0.9500^{\text{year\_diff}} && \text{for year} \leq 2030 \\
    \text{capex}_{\text{solar}} &= 700 \times 0.9500^{4} \times 0.9800^{(\text{year}-2030)} && \text{for year} > 2030
\end{align}
Baseline of 700 EUR/kWp in 2026 \citep{IEA2026}. Learning rate: $\sim$5\%/yr to 2030 ($\sim$569 EUR/kWp by 2030), then $\sim$2\%/yr thereafter \citep{DNV2025}. At 2040: $\approx$468 EUR/kWp.

\paragraph{Onshore Wind CAPEX:}
\begin{align}
    \text{capex}_{\text{wind}} &= 1200 \times 0.9775^{\text{year\_diff}} && \text{for year} \leq 2030 \\
    \text{capex}_{\text{wind}} &= 1200 \times 0.9775^{4} \times 0.9950^{(\text{year}-2030)} && \text{for year} > 2030
\end{align}
Baseline of 1{,}200 EUR/kW in 2026. Learning rate: $\sim$2.25\%/yr to 2030 ($\sim$1{,}097 EUR/kW by 2030), then $\sim$0.5\%/yr thereafter \citep{IEA2026}. At 2040: $\approx$1{,}043 EUR/kW.

\paragraph{Alkaline Electrolyser CAPEX} (default technology):
\begin{align}
    \text{capex}_{\text{elec}} &= 850 \times 0.917^{\text{year\_diff}} && \text{for year} \leq 2030 \\
    \text{capex}_{\text{elec}} &= 850 \times 0.917^{4} \times 0.976^{(\text{year}-2030)} && \text{for year} > 2030
\end{align}
Baseline: 850 EUR/kW in 2026. Learning rate: $\sim$8.3\%/yr to 2030 ($\sim$601 EUR/kW by 2030), then $\sim$2.4\%/yr thereafter \citep{DEA2026}. At 2040: $\approx$471 EUR/kW.

\paragraph{PEM Electrolyser CAPEX:}
\[ \text{capex}_{\text{elec}} = 1000 \times 0.97^{\text{year\_diff}} \]
Baseline: 1{,}000 EUR/kW in 2026. Constant learning rate of $\sim$3\%/yr \citep{DEA2026}. At 2040: $\approx$659 EUR/kW.

\paragraph{SOEC Electrolyser CAPEX:}
\[ \text{capex}_{\text{elec}} = 2200 \times 0.95^{\text{year\_diff}} \]
Baseline: 2{,}200 EUR/kW in 2026. Learning rate $\sim$5\%/yr \citep{DEA2026}. At 2040: $\approx$1{,}074 EUR/kW.

\begin{figure}[h]
    \centering
    \includegraphics[width=0.75\linewidth]{Literature Review//Figs//Section 2/S2F1.png}
    \caption{Electrolyser CAPEX trajectories \citep{DEA2026}}
    \label{fig:electrolyser-capex1}
\end{figure}

\paragraph{Electrolyser Efficiency (LHV basis):}

\begin{table}[h]
    \centering
    \caption{Electrolyser efficiency baselines and improvement rates.}
    \label{tab:elec-eff}
    \begin{tabular}{lccc}
        \toprule
        Technology & 2026 Baseline & Source \\
        \midrule
        Alkaline & 64\% & DEA, 2026\\
        PEM      & 69.2\% & DEA, 2026\\
        SOEC     & 80.8\% & DEA, 2026\\
        \bottomrule
    \end{tabular}
\end{table}

Electrolyser lifetime (in years) is derived from a technology-specific stack lifetime in operating hours divided by the annual full load hours.

\paragraph{Compression CAPEX:}
A compression CAPEX of 32 EUR/kW (2026 baseline, declining at $\sim$1.5\%/yr) is included to cover the compressor vessel and motor required to pressurise hydrogen from electrolyser outlet pressure ($\sim$30 bar) to pipeline or storage pressure (70--200 bar). This is annualised alongside the electrolyser stack via the CRF. The compression electricity demand (4 kWh/kg H\textsubscript{2}) is handled separately through the renewable electricity sizing and is not double-counted in the CAPEX.

\subsection{Location Adjustment Factors}
\label{subsec:location-adj}

Two layers of location adjustment are applied to every grid point.

\paragraph{Regional CAPEX multipliers:} Region-specific installation cost factors are applied to the global solar and wind CAPEX baselines. These reflect differences in labour costs, supply chain maturity, logistics, and permitting complexity.

\begin{table}[htbp]
    \centering
    \caption{Regional CAPEX multipliers by H2\_Region.}
    \label{tab:regional-factors}
    \small
    \begin{tabular}{lcc}
        \toprule
        H2\_Region & Solar CAPEX Factor & Wind CAPEX Factor \\
        \midrule
        EU                  & 1.30 & 1.40 \\
        Non-EU Europe       & 1.20 & 1.25 \\
        North Africa        & 0.85 & 0.85 \\
        Sub-Saharan Africa  & 1.05 & 1.05 \\
        Middle East         & 0.85 & 0.90 \\
        Central Asia        & 1.00 & 1.05 \\
        South Asia          & 0.75 & 0.85 \\
        East Asia           & 0.85 & 0.85 \\
        Southeast Asia      & 0.90 & 0.95 \\
        Oceania             & 1.10 & 1.15 \\
        North America       & 1.00 & 1.00 \\
        Latin America       & 0.90 & 0.90 \\
        Russia              & 1.10 & 1.10 \\
        \bottomrule
    \end{tabular}
\end{table}

\paragraph{Country-level WACC:} As described in \S\ref{subsec:wacc}, country-specific WACC values from \cite{Hatton2025} are applied per grid point. A grid point's effective CAPEX burden therefore depends not only on its physical location's cost multiplier, but on the financing conditions of the country it sits in. This distinguishes, for example, between Moroccan and Australian solar installations that might have similar irradiance but materially different financing costs.

The effective annualised CAPEX per unit of hydrogen production is computed as:
\[ \text{CAPEX}_{\text{annualised}} = \text{CAPEX}_{\text{global}} \times \text{location\_factor} \times \text{CRF}(\text{wacc}, \text{lifetime}) \]

where $\text{CRF}(r, n) = r / (1 - (1+r)^{-n})$.

\subsection{Electricity Demand Sizing}

The specific electrical energy and power required to produce one kilogram of hydrogen per year are:

\begin{align}
    E_{\text{spec}} &= \frac{33.3}{\text{eff}} + \text{comp\_elec} \quad [\text{kWh/kg H}_2] \\
    P_{\text{req}} &= \frac{E_{\text{spec}}}{\text{full\_load\_hours}} \quad [\text{kW}_{\text{elec}}/\text{(kg H}_2\text{/yr)}]
\end{align}

where $E_{\text{spec}}$ is the specific electrical energy demand, $\text{eff}$ is the electrolyser LHV efficiency, $\text{comp\_elec} = 4$ kWh/kg H\textsubscript{2} is the compression electricity, and full load hours = 4{,}000 hrs/yr is the assumed annual operating time of the electrolyser.

\subsection{Solar Cost Calculation}

The required solar array size and its levelised cost contribution are:
\begin{align}
    S_{\text{spec}} &= \frac{E_{\text{spec}}}{\text{Solar\_Energy\_Potential}} \quad [\text{kWp}/(\text{kg H}_2\text{/yr})] \\
    \text{CAPEX}_{\text{solar,spec}} &= S_{\text{spec}} \times \text{capex\_solar\_global} \times \text{solar\_factor} \quad [\text{EUR}/(\text{kg H}_2\text{/yr})] \\
    \text{LCOH}_{\text{solar}} &= \text{CapEx}_{\text{solar,spec}} \times \left( \text{CRF}(\text{wacc\_ren}, 25) + 0.045 \right) \quad [\text{EUR}/\text{kg H}_2]
\end{align}

where solar lifetime is 25 years, fixed O\&M is 4.5\% of CAPEX per year, and Sola Energy Potential [kWh/kWp/yr] is the location-specific photovoltaic yield from \cite{GSA}.

\subsection{Wind Cost Calculation}

Wind turbine output is derived from the location's Wind Power Density (WPD), capped at 800 W/m\textsuperscript{2} to remove data artefacts in coastal cells:

\begin{align}
    P_{\text{turbine}} &= \text{WPD}_{\text{capped}} \times 0.40 \times \pi \times (50)^2 / 10^6 \quad [\text{MW}] \\
    N_{\text{turbines}} &= P_{\text{req}} / P_{\text{turbine}} \\
    \text{Wind CAPEX} &= N_{\text{turbines}} \times 2\,[\text{MW}] \times \text{capex\_wind} \times 1000 \quad [\text{EUR}] \\
    \text{Yearly Cost Wind} &= \text{CRF}(\text{wacc\_ren}, 20) \times \text{Wind CAPEX} + 8 \times E_{\text{spec,yearly}}
\end{align}

where the Betz efficiency is set to 0.40, blade length is 50 m (representative of a modern 2 MW turbine), wind lifetime is 20 years, and variable O\&M is 8 EUR/MWh.

\paragraph{Binary electricity source selection:} Each grid point independently selects either solar or wind (whichever has a lower annualised cost) as its sole electricity source. Hybrid configurations are not modelled; this is a simplifying assumption.

\subsection{Electrolyser Cost Calculation}

The total electrolyser system CAPEX includes the stack, electrical balance-of-plant (BOP), and compression hardware:

\begin{align}
    \text{total\_capex\_h2} &= (\text{capex\_elec} + 41.6 + \text{compression\_capex}) \times P_{\text{req}} \times 1000 \quad [\text{EUR}] \\
    \text{Yearly Cost Electrolyser} &= \text{CRF}(\text{wacc\_elec}, L_{\text{elec}}) \times \text{total\_capex\_h2} \notag \\
    & \quad + 0.02 \times \text{total\_capex\_h2} + 0.07 \times \text{h2\_demand} \times 10^6
\end{align}

where \texttt{capex\_elec} is the year- and technology-adjusted electrolyser stack CAPEX [EUR/kW]; 41.6 EUR/kW is the electrical BOP (power electronics, transformers, switchgear); 2\% per year is the O\&M rate; 0.07 EUR/kg is the water OPEX; and $L_{\text{elec}}$ is the technology- and year-specific electrolyser lifetime in years.

\subsection{Total Generation LCOH}

\begin{align}
    \text{Yearly gen.\ cost} &= \min(\text{Yearly Cost Solar},\, \text{Yearly Cost Wind}) + \text{Yearly Cost Electrolyser} \\
    \text{Gen.\ cost per kg H}_2 &= \text{Yearly gen.\ cost} \,/\, (\text{h2\_demand} \times 10^6)
\end{align}

This produces the generation cost map: a global scatter of LCOH values at every grid point, with colours reflecting the spatial variation driven by renewable resource quality, CAPEX multipliers, and country-level financing costs. The cheapest locations emerge in high-irradiance, low-WACC countries such as Chile, Morocco, and Namibia.

\section{Transport Cost Model}
\label{sec:transport}

\subsection{Routing Methodology}

Transport cost parametrisation follows \cite{Collis2022}. For each grid point, the transport cost is calculated as the minimum-cost pathway from the production site to the specified European destination port across four competing carriers: ammonia (NH\textsubscript{3}), LOHC (methylcyclohexane), liquid hydrogen (LH\textsubscript{2}), and compressed H\textsubscript{2} gas (pipeline or truck). The routing involves four stages:

\begin{enumerate}
    \item \textbf{Source-to-port overland leg:} Road distance from each grid point to the nearest port is retrieved from a pre-computed cache (via the OSRM Open Source Routing Machine API using OpenStreetMap data).
    \item \textbf{Maritime leg:} Shipping distance between source and destination ports is retrieved from a pre-computed matrix of shortest navigable sea routes between all major global ports, computed via Dijkstra's algorithm on a graph network derived from global shipping lane shapefiles \citep{Collis2022}.
    \item \textbf{Carrier selection:} All four carriers are evaluated for each route; the carrier with the lowest total delivered cost is selected and its cost recorded. This carrier selection map constitutes a transport mode distribution result.
    \item \textbf{Pipeline routing:} Where corridor configurations permit pipelines and the source-to-destination distance falls within 2{,}000 km, a pipeline transport cost is evaluated as an additional alternative for H\textsubscript{2} gas and NH\textsubscript{3} carriers.
\end{enumerate}

Ship speed is assumed to be 25 km/h (typical chemical tanker transit speed), used for boil-off loss calculations.

\subsection{Transport Cost Functions}

All cost functions produce a result in EUR/kg H\textsubscript{2} delivered.

\paragraph{Ammonia (NH\textsubscript{3}):}
\begin{itemize}
    \item Conversion (source): 0.45 EUR/kg H\textsubscript{2}
    \item Export terminal: 0.11 EUR/kg H\textsubscript{2}
    \item Shipping: log-linear function of distance
    \item Pipeline (if applicable): $0.0007 \times d - 0.0697$ EUR/kg H\textsubscript{2} (where $d$ = km)
    \item Reconversion (cracking): 0.75 EUR/kg H\textsubscript{2}
    \item Boil-off gas (BOG): 0.036\%/day in transit; multiplier capped at $1.40\times$
\end{itemize}

\paragraph{LOHC (Methylcyclohexane):}
\begin{itemize}
    \item Hydrogenation at source: 0.35 EUR/kg H\textsubscript{2}
    \item Export terminal: 0.10 EUR/kg H\textsubscript{2}
    \item Shipping: log-linear function of distance
    \item Dehydrogenation: 1.10 EUR/kg H\textsubscript{2}
    \item No pipeline option; no boil-off loss
\end{itemize}

\paragraph{Liquid Hydrogen (LH\textsubscript{2}):}
\begin{itemize}
    \item Liquefaction: 0.45 EUR/kg H\textsubscript{2}
    \item Cryogenic export terminal: 0.50 EUR/kg H\textsubscript{2}
    \item Shipping: log-linear function of distance, scaled by BOG multiplier
    \item Regasification: 0.14 EUR/kg H\textsubscript{2}
    \item BOG: 0.2\%/day in transit + 5.9\% at terminal transfer; multiplier capped at $1.40\times$
\end{itemize}

\paragraph{Compressed H\textsubscript{2} Gas:}
\begin{itemize}
    \item Pipeline: $0.0004 \times d + 0.0424$ EUR/kg H\textsubscript{2} (if within distance limit)
    \item Trucking: $0.003 \times d + 0.3319$ EUR/kg H\textsubscript{2}
    \item No conversion cost; no BOG
\end{itemize}

The $1.40\times$ BOG multiplier cap prevents pathologically high shipping costs on ultra-long routes ($>$20{,}000 km). Transport cost maps and carrier mode maps are produced per corridor using these pre-computed values; the transport model is not re-run live in the dashboard.

\section{Demand Scenarios}
\label{sec:demand}

Three demand scenarios are defined for European green hydrogen imports, drawing on EHB, IEA, and EU Commission projections:

\begin{table}[htbp]
    \centering
    \caption{Demand scenarios [kt H\textsubscript{2}/yr].}
    \label{tab:demand-scenarios}
    \small
    \begin{tabular}{lccccccc}
        \toprule
        Scenario & 2025 & 2027 & 2030 & 2035 & 2040 & 2045 & 2050 \\
        \midrule
        Ambitious          & 500 & 4{,}700 & 11{,}600 & 20{,}000 & 23{,}800 & 31{,}300 & 46{,}000 \\
        Intermediate       & 400 & 2{,}700 & 5{,}900  & 9{,}500  & 12{,}100 & 16{,}700 & 25{,}000 \\
        Current Trajectory & 300 & 1{,}000 & 1{,}700  & 2{,}500  & 4{,}000  & 6{,}000  & 9{,}000  \\
        \bottomrule
    \end{tabular}
\end{table}

Intermediate years are linearly interpolated. The Ambitious scenario aligns with the REPowerEU 2030 import target. The Intermediate scenario tracks the IEA Announced Pledges scenario. The Current Trajectory scenario reflects the IEA Stated Policies scenario. At the 2040 target year, the three scenarios correspond to 23{,}800, 12{,}100, and 4{,}000 kt H\textsubscript{2}/yr respectively.

\section{Capacity Allocation and Supply Curves}
\label{sec:capacity}

\subsection{Theoretical Renewable Potential}

Country-level theoretical H\textsubscript{2} export potential is estimated from renewable resource density and land availability. A land fraction of 3\% is applied globally (representing land available for dedicated H\textsubscript{2} export projects, excluding protected areas, urban land, and agricultural land), with a higher 10\% fraction for arid countries with large uninhabited desert regions. These fractions are applied to the renewable energy yield per km\textsuperscript{2} to produce a production ceiling in kt H\textsubscript{2}/yr.

\subsection{IEA 2030 Deployment Anchors}

Near-term capacity is constrained by the IEA Hydrogen Projects Database \citep{IEA2026}. For each country, the status-weighted project capacity (per Table~\ref{tab:iea-weights}) is summed to produce a 2030 anchor value [kt H\textsubscript{2}/yr]. This anchors model projections to the current observable project pipeline rather than purely theoretical potential, and prevents the model from over-allocating supply in regions where deployment is not yet underway.

\subsection{Readiness Ramp}

The available capacity fraction for a given year is modelled as a ramp function anchored at three points:
\begin{align}
    \text{readiness}(2020) &= 0 \\
    \text{readiness}(2030) &= \text{IEA\_2030\_anchor} \;/\; \text{theoretical\_potential} \\
    \text{readiness}(2050) &= 1.0
\end{align}
Linear interpolation is applied between these anchor years, clamped to [0, 1]. A minimum floor of 1\% of theoretical potential by 2030 is enforced for all countries, preventing the complete exclusion of countries with genuinely emerging industries not yet captured in the IEA database.

The country capacity ceiling applied in the model is then:
\begin{equation}
    \text{country\_cap\_kt}(\text{year}) = \text{theoretical\_potential} \times \text{readiness}(\text{year})
\end{equation}

At 2040 this places each country roughly 60--70\% of the way along its ramp to full theoretical potential, depending on the size of its IEA 2030 anchor relative to its physical resources.

\subsection{Country Aggregation and Supply Curve Construction}

Within each corridor, grid points are grouped by ISO\_A3 (country). The representative country LCOH is a P10--P90 trimmed mean. This removes spatial interpolation artefacts and anomalous values at country edges. For example, grid cells at the coast inherit anomalously high wind power density values from the underlying Global Wind Atlas, because the 1° cell straddles land and sea. The sea portion picks up offshore wind speeds and inflates the onshore estimate.

Supply curves are then constructed as merit-order step functions:
\begin{itemize}
    \item \textbf{X-axis}: cumulative supply capacity [kt H\textsubscript{2}/yr], with each country's step width equal to its capacity
    \item \textbf{Y-axis}: representative country LCOH [EUR/kg H\textsubscript{2}]
    \item Countries are sorted cheapest-first
    \item A demand line marks the scenario target quantity
\end{itemize}

Where a country appears in multiple corridors (e.g.\ Morocco in both A and B), the corridor delivering the lower cost is used in the global optimal mix calculation. The supply curve for each demand scenario and the global optimal mix overlay are the primary analytical outputs for corridor comparison.

\section{European Hydrogen Backbone Corridor Framework}
\label{sec:corridors}

Five import gateways are defined following the European Hydrogen Backbone (EHB) network framework, each with a distinct destination port and source region. Source countries are filtered by corridor-specific allowlists:

\begin{table}[htbp]
    \centering
    \caption{EHB import gateways.}
    \label{tab:corridors}
    \small
    \begin{tabular}{p{2.5cm}p{2cm}p{4cm}}
        \toprule
        Corridor & Destination & Transport Mode \\
        \midrule
        A & Italy & Pipeline + Shipping \\
        B & Spain & Pipeline + Shipping \\
        C & Netherlands & Shipping \\
        D & Germany & Pipeline \\
        E & Austria & Pipeline \\
        EU & --- & No transport required \\
        \bottomrule
    \end{tabular}
\end{table}

EU member state grid points are excluded from all import corridors (A--E) to avoid double-counting domestic production in international supply analyses.

\section{Monte Carlo Portfolio Analysis}
\label{sec:mc-portfolio}

\subsection{Purpose and Research Questions}

The Monte Carlo portfolio analysis propagates parameter uncertainty through the full pipeline to characterise the probability distribution of the optimal supply portfolio cost. Unlike a one-at-a-time sensitivity analysis (which perturbs one parameter while holding others fixed), the MC simultaneously samples all uncertain parameters. This addresses two questions:

\begin{enumerate}
    \item What is the distribution of the capacity-weighted delivered cost of the cost-optimal supply mix at 2040? Specifically, what are the P10, P50, and P90 estimates?
    \item Which countries are always selected into the supply mix regardless of parameter values (core suppliers), which are selected only under favourable assumptions (conditional suppliers), and which rarely appear (marginal suppliers)?
\end{enumerate}

A third question, threshold feasibility, asks: under what fraction of parameter combinations does the portfolio delivered cost fall below a specified affordability threshold? And which parameter values distinguish sub-threshold runs from the full distribution?

\subsection{Parameter Sampling}

Six parameters are sampled from triangular distributions for $N = 1000$ iterations using a reproducible random number generator (seed = 42):

\begin{table}[htbp]
    \centering
    \caption{Monte Carlo parameter sampling ranges.}
    \label{tab:mc-params}
    \small
    \begin{tabular}{p{5cm}ccc{2cm}}
        \toprule
        Parameter & Low & Mode & High \\
        \midrule
        Solar CAPEX       & 325 & 465 & 605 \\
        Wind CAPEX        & 730 & 1043 & 1355 \\
        Electrolyser CAPEX & 330 & 471 & 612 \\
        Electrolyser Efficiency      & 0.65 & 0.74 & 0.82 \\
        WACC multiplier              & 0.60 & 1.00 & 1.40 \\
        Transport Cost multiplier    & 0.70 & 1.00 & 1.30 \\
        \bottomrule
    \end{tabular}
\end{table}

Triangular distributions are used because they assign the highest probability mass to the central estimate while allowing asymmetric tails. The modes reflect the IEA/IRENA/DEA central trajectory as the best estimate, with $\pm$30\% bounds capturing the spread across independent forecasters. The WACC multiplier range [0.60, 1.40] is wider than the CAPEX ranges to reflect the substantially higher uncertainty in financing conditions, which depend on macroeconomic and geopolitical factors beyond technology learning curves.

\subsection{Per-Iteration Computation}

For each of the $N = 1000$ iterations:
\begin{enumerate}
    \item Six parameters are sampled from their triangular distributions.
    \item The WACC multiplier is applied by scaling all country WACC values in place; these are restored after the iteration to preserve reproducibility.
    \item Generation costs are recomputed with CAPEX and efficiency overrides applied to the 2040 global baseline values.
    \item Total LCOH per grid point is assembled: $\text{LCOH}_{\text{total}} = \text{Gen.\ cost} + \text{transport} \times t_{\text{mult}}$, where transport costs are the pre-computed corridor values scaled by the sampled transport multiplier.
\end{enumerate}

\subsection{Country-Level Aggregation}

Within each iteration, grid points are grouped by ISO\_A3. The representative country LCOH for that iteration is the P10--P90 trimmed mean of within-cap grid point LCOHs, matching the supply curve construction methodology. Country capacity figures do not change between iterations: uncertainty is applied only to cost parameters, not to deployment constraints.

\subsection{Greedy Portfolio Construction}

For each iteration, a greedy cost-optimal supply portfolio is built:
\begin{enumerate}
    \item Countries are sorted in ascending order of trimmed-mean LCOH.
    \item Starting from the cheapest country, 
    demand is allocated up to the minimum of remaining demand or country capacity.
    \item Allocation continues until the total demand target is met or all countries are exhausted.
\end{enumerate}

Outputs per iteration: capacity-weighted average portfolio LCOH, number of countries selected, total demand met, regional Herfindahl-Hirschman Index ($\text{HHI} = \sum_r s_r^2$ where $s_r$ is the regional share of allocated supply), and the set of selected ISO\_A3 codes.

\subsection{Country Classification}

Across all 500 iterations, each country's selection frequency is computed. Countries are classified into three tiers:
\begin{itemize}
    \item \textbf{Core suppliers:} selected in $>$90\% of simulations --- appear in the optimal mix regardless of parameter uncertainty
    \item \textbf{Conditional suppliers:} selected in 50--90\% of simulations --- typically selected but sensitive to assumptions
    \item \textbf{Marginal suppliers:} selected in $<$50\% of simulations --- only appear under favourable conditions
\end{itemize}

\subsection{Threshold Feasibility Analysis}

For a specified price threshold (EUR/kg H\textsubscript{2}), the fraction of simulations achieving LCOH less than or equal to the threshold is reported. For this analysis, a threshold of €3/kg is chosen as. Histograms are then used to compare the distribution of each parameter value across all runs versus the subset of runs that achieve the threshold. Parameters whose threshold distribution is visibly shifted relative to the full distribution are necessary conditions for affordability. Their values must be toward the optimistic end of the range for the target to be achievable. Parameters whose distributions are unchanged across the two groups have little influence on whether the threshold is met.

\section{Multi-Objective Portfolio Analysis}
\label{sec:mc-multi}

\subsection{Purpose}

This analysis extends the cost-optimal portfolio MC to test whether imposing non-cost objectives (supply security, water availability) materially changes the portfolio cost distribution or the identity of selected suppliers. It applies three objective weightings to the same $N = 1000$ MC iterations, using identical parameter samples. The LCOH is computed once per iteration; the portfolio is then selected three times under different objective functions. This design isolates the effect of the objective weighting from parameter uncertainty.

\subsection{Objective Scenarios}

\begin{table}[htbp]
    \centering
    \caption{Multi-objective portfolio scenarios.}
    \label{tab:mc-objectives}
    \small
    \begin{tabular}{lccc}
        \toprule
        Scenario & $w_{\text{cost}}$ & $w_{\text{sec}}$ & $w_{\text{water}}$ \\
        \midrule
        Cost-Optimal        & 100 & 0  & 0  \\
        70:30 Cost:Security & 70  & 30 & 0  \\
        50:30:20 C:S:W      & 50  & 30 & 20 \\
        \bottomrule
    \end{tabular}
\end{table}

A fixed diversification weight $w_{\text{dep}} = 25$ is applied in all three scenarios, imposing a concentration cap that prevents any single country from supplying more than 76.25\% of total demand in the first-pass allocation.

\subsection{Composite Score Normalisation}

For each country in each iteration, three normalised scores are computed:

\begin{align}
    \text{cost\_norm}     &= \text{clip}\!\left( \frac{\text{lcoh} - \text{lcoh}_{\min}}{\text{lcoh}_{p95} - \text{lcoh}_{\min}},\, 0,\, 1 \right) \\
    \text{security\_norm} &= \text{clip}(\text{cpi\_score} / 100,\, 0,\, 1) \\
    \text{water\_norm}    &= \text{clip}(\text{bws\_score} / 5.0,\, 0,\, 1)
\end{align}

where $\text{lcoh}_{p95}$ is the 95th percentile LCOH across countries which is used instead of the maximum to limit the influence of extreme outliers. CPI and BWS values are loaded once at run start and are static across iterations; only cost changes between iterations.

The composite score (lower = preferred) is:
\begin{align}
    \text{composite} &= \frac{w_{\text{cost}}}{w_{\text{sum}}} \times \text{cost\_norm} \notag \\
                     &+ \frac{w_{\text{sec}}}{w_{\text{sum}}} \times (1 - \text{security\_norm}) \notag \\
                     &+ \frac{w_{\text{water}}}{w_{\text{sum}}} \times \text{water\_norm}
\end{align}

where $w_{\text{sum}} = w_{\text{cost}} + w_{\text{sec}} + w_{\text{water}}$. The directionality is: cost and water stress are penalised directly (higher = worse); security is penalised inversely (higher CPI = more secure = preferred, hence the $(1 - \text{security\_norm})$ term).

\subsection{Concentration Cap and Two-Pass Allocation}

The diversification weight controls a per-country capacity cap:
\[ \text{max\_per\_country\_kt} = \text{H2\_demand} \times \max(0.05,\, 1.0 - 0.95 \times w_{\text{dep}}/100) \]

For $w_{\text{dep}} = 25$: $\text{max\_per\_country\_kt} = 0.7625 \times \text{H2\_demand}$. Countries are allocated in ascending composite score order up to this cap in a first pass. If demand remains unmet, a second pass relaxes the cap and fills the remainder from whichever countries have residual capacity, again in composite score order.

\subsection{Outputs}

For each objective scenario across each corridor--demand combination:
\begin{itemize}
    \item \textbf{Overlaid cost distribution:} histogram of portfolio \texttt{weighted\_lcoh} across $N = 500$ simulations with P10/P50/P90 markers for all three weightings on the same axes, enabling direct comparison of cost distributions under different objectives
    \item \textbf{Country selection frequency:} horizontal bar chart of percentage of simulations in which each country was selected, per objective weighting, with core/conditional/marginal tier lines
    \item \textbf{Threshold feasibility:} per objective, the fraction of simulations achieving $\leq$ threshold and per-parameter overlaid histograms (all runs vs sub-threshold runs)
\end{itemize}

\section{Strategic Multi-Criteria Dispatch and Scenario Matrix}
\label{sec:strategic-dispatch}

\subsection{Purpose}

The strategic dispatch analysis applies multi-criteria objective functions to the deterministic (point-estimate) 2040 costs rather than sampling from uncertainty distributions. This produces the 3$\times$3 scenario matrix results: world maps and radar charts for each combination of demand scenario and objective weighting. The dispatch logic is identical in structure to the MC portfolio objectives but uses fixed cost values, making results directly comparable across demand scenarios without the noise of parameter sampling.

\subsection{Country-Level Aggregation}

For each corridor, grid points are grouped by country and representative costs computed using the P10--P90 trimmed mean, consistent with supply curve construction. Where a country appears in more than one corridor, the corridor delivering the lowest representative LCOH is retained, and the country is assigned to that corridor.

CPI scores and BWS scores are merged to this country table from the auxiliary datasets (\S\ref{subsec:wacc}).

\subsection{Composite Score and Dispatch Algorithm}

The normalisation and composite score formulation are identical to \S\ref{sec:mc-multi}, applied to the deterministic LCOH values. The three objective weightings evaluated are the same as the weightings in Table 3.7.

Countries are ranked by composite score (ascending) and allocated in greedy order, subject to the same two-pass concentration cap (\S\ref{sec:mc-multi}) with $w_{\text{dep}} = 25$.

\subsection{Portfolio Key Performance Indicators}

From each dispatch result, five KPIs are derived:
\begin{itemize}
    \item \textbf{Delivered cost:} mean LCOH across allocated countries [EUR/kg H\textsubscript{2}]
    \item \textbf{Weighted security:} mean CPI score
    \item \textbf{Supply diversity:} Herfindahl-Hirschman Index by H2\_Region, $\text{HHI} = \sum_r s_r^2$ where $s_r$ is each region's share of total allocated supply
    \item \textbf{Weighted water stress:} mean BWS score
    \item \textbf{Total Carbon Avoided:} the quantity of CO\textsubscript{2} avoided relative to a grey hydrogen reference [tCO\textsubscript{2}]
\end{itemize}

\subsection{Radar Normalisation}

Each KPI is mapped to a 0--1 score (higher = better) for radar chart visualisation:

\begin{align}
    \text{radar\_cost}            &= \text{clip}\!\left(1 - \frac{\text{delivered\_cost} - 0.5}{5.0},\, 0,\, 1\right) \\
    \text{radar\_security}        &= \text{clip}\left(\frac{\text{weighted\_CPI}}{100},\, 0,\, 1\right) \\
    \text{radar\_diversification} &= \text{clip}(1 - \text{HHI},\, 0,\, 1) \\
    \text{radar\_water}           &= \text{clip}\!\left(1 - \frac{\text{weighted\_BWS}}{5.0},\, 0,\, 1\right) \\
    \text{radar\_carbon}          &= \text{clip}\!\left(1 - \frac{\text{carbon\_avoided\_cost}}{500},\, 0,\, 1\right)
\end{align}

The cost axis is normalised on a [3, 5] EUR/kg range, placing a delivered cost of 3 EUR/kg at a score of 1.0 and a cost above 5 EUR/kg at 0. The carbon axis uses a 500 EUR/tCO\textsubscript{2} reference abatement cost, consistent with ambitious decarbonisation scenarios.

\subsection{3$\times$3 Scenario Matrix Output}

The full results are presented as two 3$\times$3 matrices:
\begin{itemize}
    \item \textbf{Rows}: three demand scenarios (Ambitious 23{,}800 kt/yr, Intermediate 12{,}100 kt/yr, Current Trajectory 4{,}000 kt/yr)
    \item \textbf{Columns}: three objective weightings (Cost Only, Cost + Security, Cost + Security + Water)
\end{itemize}

Each cell of the \textbf{map matrix} shows a world map with supply countries marked as circles: circle size is proportional to the allocated quantity [kt H\textsubscript{2}/yr], and circle colour indicates the representative LCOH on a shared colour scale (1.5--4.5 EUR/kg). This reveals both which countries are selected and at what cost under each scenario-objective combination.

Each cell of the \textbf{radar matrix} shows the five-axis spider chart for the portfolio selected in that cell. Comparing across columns shows the trade-off between cost performance and non-cost metrics as objective weights shift; comparing across rows shows how increasing demand forces the portfolio into progressively more expensive or less secure supply regions.

\section{Interactive Dashboard}
\label{sec:dashboard}

\subsection{Overview}

The model outputs are presented through an interactive Streamlit dashboard. The dashboard reads pre-computed result files and allows the user to explore the full output space without re-running the pipeline. It is organised into a persistent sidebar for global scenario controls and nine analytical tabs.

\subsection{Sidebar Controls}

The sidebar governs the global state of the dashboard and applies to all tabs unless a tab provides its own override:

\begin{itemize}
    \item \textbf{Demand scenario:} selects one of the three demand scenarios (Ambitious, Intermediate, Current Trajectory).
    \item \textbf{Model year:} selects the target year for the scenario; the corresponding demand quantity is displayed as a metric.
    \item \textbf{Corridor toggles:} five checkboxes enable or disable individual EHB import corridors (A--E) and EU domestic production independently.
    \item \textbf{Apply country capacity limits:} toggle that restricts cost calculations to grid points within each country's readiness-adjusted capacity ceiling.
    \item \textbf{Electrolyser technology:} radio selector between Alkaline, PEM, and SOEC, which sets the default CAPEX and efficiency assumptions.
    \item \textbf{CAPEX sliders:} three sliders for Solar CAPEX (100--1{,}200 EUR/kWp), Wind CAPEX (400--2{,}000 EUR/kW), and Electrolyser CAPEX (100--1{,}500 EUR/kW), pre-populated from the technology cost curves at the selected year and electrolyser type. A reset button restores all three to their computed defaults.
    \item \textbf{Transport cost adjustment:} slider from $-50\%$ to $+50\%$ applied as a uniform multiplier to all pre-computed transport costs, enabling sensitivity testing without re-running the routing model.
    \item \textbf{RED III limit:} slider for the renewable fuel carbon intensity threshold (0.5--10.0 kgCO\textsubscript{2}eq/kgH\textsubscript{2}), with a default of 3.38 kgCO\textsubscript{2}eq/kgH\textsubscript{2} matching the current EU delegated regulation.
\end{itemize}

\subsection{Dashboard Tabs}

\paragraph{Tab 1 — Supply Curve:} Displays the merit-order supply curve for the selected corridors, scenario, and year. Each country appears as a step, sorted by representative LCOH, with cumulative capacity on the x-axis and delivered cost on the y-axis. A horizontal demand line marks the scenario target.

\paragraph{Tab 2 — Cost Breakdown:} Shows a stacked bar chart decomposing total delivered cost into generation cost components (solar/wind, electrolyser) and transport cost for each country in the selected portfolio. Enables identification of whether a given supplier's competitiveness is driven by cheap renewables, low financing costs, or short transport distances.

\paragraph{Tab 3 — Source Maps:} Renders four global maps: generation cost, transport cost, total delivered cost, and optimal transport mode. A selectbox allows the user to change the destination import port (Huelva, Civitavecchia, Vienna, Rotterdam, or Hamburg), which updates the transport and total cost maps while leaving the generation cost map unchanged.

\paragraph{Tab 4 — Transport Modes:} Displays a pie chart and summary table of the carrier mode distribution (ammonia, LOHC, liquid hydrogen, compressed gas) across all grid points in the selected corridor set. No tab-specific controls; reflects the sidebar CAPEX and transport multiplier settings.

\paragraph{Tab 5 — Country Caps:} Visualises country-level capacity ceilings as a bar chart. A year slider (ranging over the full model period) allows the user to inspect how the readiness ramp evolves over time independently of the main scenario year selected in the sidebar.

\paragraph{Tab 6 — Flow Map:} Renders a Sankey-style flow map showing the routing of supply from source countries through corridors to European destination ports. A slider controls the number of top source countries displayed per corridor (5--30).

\paragraph{Tab 7 — Assumptions:} Diagnostic tab displaying the technology cost curve trajectories used in the model. A radio selector switches between Alkaline, PEM, and SOEC to inspect the CAPEX and efficiency profiles of each electrolyser technology independently of the sidebar selection.

\paragraph{Tab 8 — Strategic Dispatch:} Implements the deterministic multi-criteria dispatch described in \S\ref{sec:strategic-dispatch} with live user-configurable objective weights. Controls include:
\begin{itemize}
    \item Four objective weight sliders (Cost, Security, Water Access, Diversification), each 0--100; normalised internally so only relative magnitudes matter.
    \item A diversification mode selectbox (by Country, Region, or Corridor) that determines the entity over which the concentration cap is applied.
    \item A checkbox to exempt EU domestic production from the diversification constraint.
    \item A capacity limits toggle (independent of the sidebar toggle).
    \item Two number inputs for a grey hydrogen cost and emissions reference, used to compute the carbon avoided KPI displayed in the radar chart.
    \item A PDF report generator that compiles the dispatch map, radar chart, and portfolio table into a downloadable report.
\end{itemize}

\paragraph{Tab 9 — H\textsubscript{2} Projects Pipeline:} Displays the IEA Hydrogen Projects Database filtered to green hydrogen projects. A multiselect filter allows the user to include or exclude projects by status (Operational, FID/Construction, Feasibility study, Concept, etc.), and a slider controls the number of top countries shown in the capacity bar chart.

\section{Key Assumptions Summary}
\label{sec:assumptions}

\begin{table}[htbp]
    \centering
    \caption{Key model assumptions summary.}
    \label{tab:assumptions}
    \small
    \begin{tabular}{p{4cm}p{4cm}p{6cm}}
        \toprule
        Assumption & Value / Approach & Justification \\
        \midrule
        Grid resolution & $\sim$1$\degree$ global, 5{,}970 onshore + 16 offshore & Computational tractability; captures country-level spatial variation \\
        Target year & 2040 & Aligns with EU 2040 climate target milestones \\
        Electricity source & Binary (solar OR wind per location) & Simplification; no hybrid optimisation \\
        Full load hours & 6{,}000 hrs/yr & Consistent capacity factor for coupled renewable-electrolyser systems \\
        BOG multiplier cap & 1.40$\times$ & Prevents artefactual extreme costs at $>$20{,}000 km routes \\
        Country aggregation & P10--P90 trimmed mean & Removes spatial outliers at country boundaries \\
        Demand allocation denominator & $n = 5{,}986$ global points & Consistent per-point volume across geographically restricted corridors \\
        Capacity readiness & IEA 2030 anchor + linear ramp to 100\% by 2050 & Anchors near-term deployment to current project pipeline \\
        Minimum readiness floor & 1\% of theoretical by 2030 & Prevents complete exclusion of emerging producers \\
        MC iterations & $N = 500$ & Sufficient for stable P10/P50/P90 convergence \\
        MC parameter distributions & Triangular (mode = 1.0 for multipliers) & Highest probability at central IEA/IRENA estimate \\
        Diversification weight & $w_{\text{dep}} = 25$ (all portfolio analyses) & Limits single-country supply share to 76.25\% \\
        Missing CPI/BWS & Imputed at sample median & Avoids excluding countries without data whilst preserving normalisation \\
        RNG seed & 42 (all MC analyses) & Full reproducibility \\
        \bottomrule
    \end{tabular}
\end{table}
